\section{Cosmology: The Global $\sigma > 0$ Phase}
\label{sec:cosmology}

\subsection{Shift in the Question}

Section~\ref{sec:black_hole} identified the black hole interior as a local
non-equilibrium closure phase characterized by $\sigma > 0$.
This raises the natural global question:
\begin{equation}
  \boxed{
    \text{What happens when the } \sigma > 0
    \text{ phase extends across the entire universe?}
  }
\end{equation}
Within the VED framework, the answer is direct:
\begin{equation}
  \boxed{\text{Cosmic expansion emerges as a global } \sigma > 0 \text{ phase.}}
\end{equation}
Thus cosmology is interpreted not as an initial condition,
but as a phase of the closure dynamics.

\subsection{The Fundamental Equation}

From Vol.~1, the continuity equation for log density is
$\partial_t \rho + \nabla \cdot J = \sigma$.
Under the homogeneous and isotropic approximation,
$\nabla \cdot J \approx 3H\rho$, which yields:
\begin{equation}
  \label{eq:continuity_cosmo}
  \boxed{\dot{\rho} + 3H\rho = \sigma.}
\end{equation}
The first term represents dilution due to expansion;
the second represents generation of structure.

\subsection{The Equation in $\tau$-Time}

Using the time definition from Vol.~1, $d\tau/dt = \rho$,
the equation becomes:
\begin{equation}
  \label{eq:tau_evolution}
  \boxed{\frac{d\rho}{d\tau} + 3H_\tau \rho = \sigma(\rho).}
\end{equation}
The scale factor is defined structurally as
$a(\tau) \sim \langle d_{ij}(\tau)\rangle$.
Since the expansion rate is determined by density, $H_\tau \sim \rho$, giving:
\begin{equation}
  H^2 \propto \rho,
\end{equation}
which is structurally analogous to the Friedmann equation.

\subsection{Final Form}

Substituting $H_\tau \sim \rho$, the evolution equation becomes:
\begin{equation}
  \label{eq:VED_cosmo}
  \boxed{\frac{d\rho}{d\tau} + 3\gamma \rho^2 = \sigma(\rho).}
\end{equation}
Here $3\gamma\rho^2$ represents dilution due to expansion
and $\sigma(\rho)$ represents generation of logs.
Cosmology therefore emerges as competition between generation and dilution.

\subsection{The Form of $\sigma$}

From the self-referential structure of Vol.~1,
the minimal effective form is:
\begin{equation}
  \label{eq:sigma_form}
  \boxed{\sigma(\rho) = \xi \rho \left(1 - \frac{\rho}{\rho_{\max}}\right).}
\end{equation}
This corresponds to growth at low density and saturation at high density.
The logistic form arises naturally from self-referential selection.

\subsection{Fixed Point and de Sitter Phase}

The fixed-point solution satisfies $3\gamma\rho_*^2 = \sigma(\rho_*)$, giving:
\begin{equation}
  \rho_* = \frac{\xi}{\xi/\rho_{\max} + 3\gamma}.
\end{equation}
At the fixed point, $H_* = \gamma\rho_* = \text{const}$, and therefore:
\begin{equation}
  \label{eq:deSitter}
  \boxed{a(\tau) \sim e^{H_*\tau}.}
\end{equation}
\begin{equation}
  \boxed{\text{de Sitter-type accelerated expansion.}}
\end{equation}

\subsection{Origin of the Cosmological Constant}

The effective cosmological constant is:
\begin{equation}
  \label{eq:lambda_eff}
  \boxed{\Lambda_{\text{eff}} \sim \gamma^2\rho_*^2.}
\end{equation}
The cosmological constant is therefore not fundamental,
but determined by closure structure.

\subsection{The Inversion}

Conventional cosmology: expansion = initial condition.

VED:
\begin{equation}
  \boxed{\text{Cosmic expansion = phase selection.}}
\end{equation}
If $\sigma = 0$, the universe is static.
If $\sigma > 0$, the universe expands.

\subsection{Unification with Black Holes}

The black hole interior satisfies $\sigma > 0$;
the expanding universe also satisfies $\sigma > 0$.
\begin{equation}
  \boxed{
    \text{Black hole interior and cosmic expansion are the same phase
    at different scales.}
  }
\end{equation}

\subsection{Conclusion}

Cosmology emerges as global closure dynamics.
\begin{equation}
  \boxed{\text{The universe is continuously generated.}}
\end{equation}
\begin{equation}
  \boxed{\text{Expansion is a consequence of } \sigma > 0.}
\end{equation}
The large-scale dynamics of the universe follow directly from closure structure.
